% !TEX program = xelatex
\documentclass[12pt, a4paper]{article}

% --- Paquetes ---
\usepackage[spanish]{babel}
\usepackage[utf8]{inputenc}
\usepackage[T1]{fontenc}
\usepackage{textcomp}
\usepackage{float}
\usepackage{enumitem}
\usepackage{longtable}
\usepackage{array}
\usepackage{booktabs}
\usepackage{tabularx}
\usepackage{xltabular}
\usepackage{colortbl}
\usepackage{multirow}
\usepackage{setspace}
\usepackage{titlesec}
\usepackage{caption}
\usepackage{graphicx}

% --- Tipografia: Lux Sans (xelatex obligatorio) ---
\usepackage{fontspec}
\setmainfont{LuxSans-Regular}[
  Path = fonts/,
  Extension = .ttf,
  BoldFont = LuxSans-Bold,
  ItalicFont = LuxSans-Regular,
  ItalicFeatures = {FakeSlant=0.2},
  BoldItalicFont = LuxSans-Bold,
  BoldItalicFeatures = {FakeSlant=0.2},
]

% --- Margenes ---
\usepackage[a4paper,top=3cm,bottom=3cm,left=3cm,right=3cm]{geometry}

% --- Hiperlinks ---
\usepackage[colorlinks=true, linkcolor=black, citecolor=blue, urlcolor=blue]{hyperref}

% --- Colores ---
\usepackage[table]{xcolor}
\definecolor{azulPrincipal}{RGB}{28, 61, 90}
\definecolor{grisClaro}{RGB}{230, 230, 230}

% --- Legibilidad ---
\setstretch{1.1}
\setlength{\parskip}{6pt}
\setlength{\parindent}{0pt}
\captionsetup{justification=centering}
\titlespacing{\section}{0pt}{16pt}{0pt}
\titlespacing{\subsection}{0pt}{12pt}{0pt}
\titlespacing{\subsubsection}{0pt}{10pt}{0pt}

% ================================================================
\begin{document}

% ================================================================
\section{Etapas}


\vspace{0.2cm}
{\renewcommand{\arraystretch}{1.4}
\setlength{\aboverulesep}{0pt}
\setlength{\belowrulesep}{0pt}
\begin{xltabular}{\textwidth}{>{\raggedright\arraybackslash}p{3.2cm} >{\raggedright\arraybackslash}X >{\raggedright\arraybackslash}p{2.4cm}}
\toprule
\rowcolor{grisClaro}
\textbf{Etapa} & \textbf{Qué se hizo} & \textbf{Archivo} \\
\hline
\endhead
1. Relevamiento topográfico & Se importaron los 35 puntos con cota, se armó la superficie del
terreno (TIN) con curvas de nivel y se dibujó el perímetro como borde de la superficie. &
\texttt{etapa1.dwg} \\
\hline
2. Proyecto de loteo & La IA diseñó la subdivisión cumpliendo la normativa: 3 calles nuevas,
4 manzanas, 76 lotes y una plaza del 10,46 \%. Después la pasó a Civil 3D como 78 parcelas
reales. & \texttt{etapa2.dwg} \\
\hline
3. Obra lineal de la calle A & Sobre la calle principal: eje, perfil del terreno, rasante al
$-$1,2 \%, sección tipo (calzada, cordón cuneta, vereda y taludes), corredor y cálculo del
movimiento de suelo. & \texttt{etapa3.dwg} \\
\hline
4. Desagüe pluvial & Red de conductos por gravedad hacia el punto bajo del terreno: 7 cámaras,
7 sumideros, 1 descarga y 14 conductos de 0,40 a 0,80 m, dimensionados por el método
racional. & \texttt{etapa4.dwg} \\
\hline
5. Alumbrado y arbolado & 23 columnas de alumbrado alternadas entre veredas, con separación
máxima de 30 m, y 140 árboles cada 8 m, respetando esquinas, sumideros y cámaras. &
\texttt{etapa5.dwg} \\
\bottomrule
\end{xltabular}}

% ================================================================
\section{Archivos}

\subsection{Datos de entrada}

\begin{itemize}[itemsep=4pt]
  \item \textbf{\texttt{puntos-relevamiento.csv}}: los 35 puntos del relevamiento, en una
  grilla cada 40 m, en formato PNEZD (número, norte, este, cota y descripción). El terreno es
  un plano que baja hacia el este y hacia el norte; el punto más bajo está en la esquina
  noreste.
  \item \textbf{\texttt{terreno-poligono.csv}}: las coordenadas de los 4 vértices del
  perímetro del terreno.
  \item \textbf{\texttt{normativa.md}}: la normativa municipal inventada que tenía que cumplir
  el loteo: anchos de calle y vereda, largo máximo de manzana, lote mínimo, espacio verde y
  criterios para alumbrado, arbolado y desagüe.
\end{itemize}

\subsection{Dibujos de Civil 3D}

Cada dibujo contiene todo lo de la etapa anterior más lo nuevo de su etapa. Todos están en
metros, creados con la plantilla métrica de Civil 3D.

\begin{itemize}[itemsep=8pt]
  \item \textbf{\texttt{etapa1.dwg}} --- el terreno.
  \begin{itemize}[itemsep=2pt]
    \item 35 puntos COGO en el grupo \textit{Relevamiento}, en una grilla cada 40 m, con
    cotas entre 96,16 y 100,00 m.
    \item Superficie TIN \textit{Terreno Natural} de 48 triángulos, con curvas de nivel
    secundarias cada 0,25 m y principales cada 1 m (estilo creado para la prueba).
    \item Perímetro del terreno como polilínea cerrada de 240 × 160 m en la capa
    \texttt{PERIMETRO}, aplicado como borde exterior de la superficie. Área:
    38.400~m\textsuperscript{2}.
  \end{itemize}

  \item \textbf{\texttt{etapa2.dwg}} --- el loteo.
  \begin{itemize}[itemsep=2pt]
    \item Emplazamiento \textit{Loteo} con 78 parcelas de Civil 3D.
    \item Manzanas M1 y M2 (112,5 × 60 m): 22 lotes cada una, 20 de 10 × 30 m
    (300~m\textsuperscript{2}) y 2 de esquina de 371~m\textsuperscript{2}.
    \item Manzana M3 (112,5 × 70 m): 10 lotes de 11 × 35 m (385~m\textsuperscript{2}) y el
    espacio verde \texttt{EV-01} de 4.017~m\textsuperscript{2} (10,46 \%).
    \item Manzana M4 (112,5 × 70 m): 22 lotes, 20 de 10 × 35 m (350~m\textsuperscript{2}) y
    2 de esquina de 433,5~m\textsuperscript{2}.
    \item Una sola parcela \texttt{CALLES} de 9.182~m\textsuperscript{2}, con las tres
    calles de 15 m y las ochavas de 4 m.
    \item Los lotes se nombran por manzana y número (M1-L01, M1-L02, etc.): primero la fila
    sur, de oeste a este, y después la fila norte.
    \item Ejes de las calles como polilíneas en la capa \texttt{EJES}: la calle A es N-S y
    conecta con la calle existente del sur; B y C son E-O.
  \end{itemize}

  \item \textbf{\texttt{etapa3.dwg}} --- la calle A como obra lineal.
  \begin{itemize}[itemsep=2pt]
    \item Alineamiento \textit{Calle A} de 152,5 m, en línea recta de sur a norte.
    \item Perfil del terreno \textit{Calle A - TN} y rasante \textit{Calle A - Rasante}, de
    una sola pendiente de $-$1,20 \% (de 99,04 a 97,21 m), con su vista de perfil.
    \item Ensamblaje \textit{Calle A - Seccion tipo}, simétrico, con 4 subensamblajes por
    lado: calzada de 3,00 m al $-$2 \% (\texttt{BasicLane}), cordón cuneta con cuneta de
    0,50 m (\texttt{UrbanCurbGutterGeneral}), vereda de 3,80 m al +2 \%
    (\texttt{UrbanSidewalk}) y talud 1:2 hasta el terreno (\texttt{DaylightGeneral}).
    Total: 15,00 m entre líneas municipales.
    \item Corredor \textit{Calle A}, con secciones cada 5 m y el terreno natural como
    objetivo de los taludes.
    \item Superficies del corredor (rasante terminada y subrasante), usadas para el
    movimiento de suelo: 37,7~m\textsuperscript{3} de desmonte y 149,4~m\textsuperscript{3}
    de terraplén contra la rasante terminada.
  \end{itemize}

  \item \textbf{\texttt{etapa4.dwg}} --- el desagüe pluvial.
  \begin{itemize}[itemsep=2pt]
    \item Red de tuberías por gravedad \textit{Pluvial Loteo}, de hormigón, con el terreno
    natural como superficie de referencia.
    \item Recorrido: calle B (tramo este), calle A hacia el norte y calle C (tramo este),
    hasta la descarga libre en la esquina noreste, punto bajo del terreno.
    \item 7 cámaras de inspección (CI-1 a CI-7), 7 sumideros (S1 a S7) y 1 cabecera de
    descarga (\texttt{DESC}).
    \item 14 conductos: 7 del colector principal (P-01 a P-07), con diámetros de 0,50 a
    0,80 m y pendientes de 0,3 a 1 \%, y 7 conexiones de sumideros (L-S1 a L-S7) de 0,40 m.
    \item Dimensionada por el método racional: caudal de diseño de
    0,64~m\textsuperscript{3}/s en la descarga y tapada mínima de 1,00 m en todos los conductos.
  \end{itemize}

  \item \textbf{\texttt{etapa5.dwg}} --- alumbrado y arbolado (resultado final).
  \begin{itemize}[itemsep=2pt]
    \item 23 columnas de alumbrado: bloque \texttt{COLUMNA\_AP} (círculo de 0,6 m con cruz)
    en la capa \texttt{ALUMBRADO}, a 0,5 m de la cara del cordón. Van alternadas entre
    veredas, con columna en cada cruce y vanos de 30 m como máximo: 7 en la calle A, 9 en la
    B y 9 en la C, contando una sola vez las de los cruces.
    \item 140 árboles: bloque \texttt{ARBOL\_2M} (círculo de 2 m) en la capa
    \texttt{ARBOLADO}, a 1,5 m de la cara del cordón y cada 8 m en las dos veredas, a 5 m
    como mínimo de las esquinas: 26 en la calle A, 56 en la B y 58 en la C.
  \end{itemize}
\end{itemize}


\end{document}
